\title{The Era of 1-bit LLMs: All Large Language Models are in 1.58 Bits}

\begin{abstract}
Recent advancements, exemplified by BitNet~\cite{bitnet}, are ushering in a transformative phase for 1-bit Large Language Models (LLMs). In this study, we propose a ternary LLM architecture called \textbf{\text{BitNet b1.58}}, where each parameter (or weight) of the LLM takes values from \{-1, 0, 1\}. This architecture achieves comparable perplexity and downstream task accuracy to full-precision Transformer LLMs (FP16 or BF16) with the equivalent parameter count and number of training tokens, yet offers substantial improvements in latency, memory footprint, throughput, and energy efficiency. Notably, the 1.58-bit LLM establishes a novel scaling law and provides new guidelines for developing high-performing, cost-efficient future LLMs. Moreover, it paves the way for new computational paradigms and facilitates the design of specialized hardware dedicated to 1-bit LLMs.
\end{abstract}

\section{The Era of 1-bit LLMs}

The recent surge in the scale and capability of Large Language Models (LLMs) has resulted in exceptional progress across diverse natural language processing benchmarks. However, the escalating model sizes have introduced significant hurdles for model deployment and highlighted environmental and economic sustainability issues due to their intensive energy demands.

To mitigate these concerns, post-training quantization has emerged as a strategy to generate low-bit models for inference~\cite{smoothquant, gptq, quip, quip_sharp}. This methodology compresses weight and activation precision, thereby decreasing the memory usage and computational load required for LLMs. The industry’s focus has shifted from 16-bit down to even 4-bit implementations~\cite{gptq, awq}. Despite its prevalence in industrial applications, post-training quantization has inherent limitations and is not the most optimal approach.

Recent advancements in 1-bit model designs, exemplified by BitNet~\cite{bitnet}, chart a promising course toward lowering the operational costs of LLMs without compromising their effectiveness. Conventional LLMs utilize 16-bit floating-point formats (such as FP16 or BF16), with matrix multiplication comprising the vast majority of their computations. This makes floating-point additions and multiplications the principal contributors to computational expense. In contrast, BitNet's matrix multiplication relies exclusively on integer addition, significantly reducing energy consumption for LLM workloads. Given that power often represents a fundamental constraint on computational performance in modern chips, such energy efficiency also translates into enhanced processing speed.

Beyond the raw computation, the transfer of model parameters from off-chip DRAM into the accelerator's on-chip memory (such as SRAM) constitutes a nontrivial cost during inference. While some strategies have explored expanding SRAM to improve throughput, the associated expenses are much steeper than for DRAM. Compared with their full-precision counterparts, 1-bit LLMs require substantially less memory in both capacity and bandwidth, drastically curtailing the overhead and latency of loading model weights from DRAM and thereby supporting faster, more cost-effective inference.

In this study, we present a noteworthy extension to the family of 1-bit LLMs, designated \textbf{\text{BitNet b1.58}}. In this architecture, every model parameter assumes one of three possible values from \{-1, 0, 1\}, as opposed to the original binary setting. By augmenting the value set with an explicit 0, BitNet b1.58 attains an effective 1.58 bits per parameter in the binary encoding. All the key advantages of BitNet are retained: the novel computation scheme largely obviates the need for multiplication in matrix operations and can be aggressively optimized; energy consumption remains equivalent to that of the 1-bit variant; and memory, bandwidth, and latency performance far outstrip that of FP16-based LLMs. Notably, \text{BitNet b1.58} provides two further benefits. First, it exhibits enhanced modeling expressivity, thanks to the capability for explicit feature suppression via zero-valued weights, substantially boosting 1-bit LLM performance. Second, experimental results demonstrate that, beginning at the 3B parameter scale and under equivalent training configurations (e.g., model size, number of training tokens), \text{BitNet b1.58} can achieve parity with full-precision (FP16) models in both perplexity and downstream task performance.

\section{BitNet b1.58}

\text{BitNet b1.58} extends the BitNet framework, a Transformer variant in which the standard \emph{nn.Linear} modules are replaced by \emph{BitLinear} layers. This model is trained from initialization, utilizing 1.58-bit quantized weights alongside activations quantized to 8 bits. In relation to the original BitNet, several adjustments are made, which are outlined below.

\paragraph{Quantization Function.} To limit the weights strictly to the values -1, 0, or +1, an \emph{absmean} quantization approach is applied. Here, the weight matrix is divided by its mean absolute value prior to rounding each entry to the nearest member of the set $\{-1, 0, +1\}$:

\begin{equation}
    \widetilde{W} = \mathrm{RoundClip}(\frac{W}{\gamma + \epsilon}, -1, 1),
\end{equation}

\begin{equation}
    \mathrm{RoundClip}(x, a, b) = \max(a, \min(b, \mathrm{round}(x))),
\end{equation}

\begin{equation}
    \gamma = \frac{1}{nm}\sum_{ij} |W_{ij}|.
\end{equation}

The quantization strategy for activations is kept consistent with BitNet; however, unlike the prior implementation, activations are not rescaled to $[0, Q_b]$ ahead of nonlinear operations. Instead, for each token, activations are normalized to the range $[-Q_b, Q_b]$, thereby eliminating the requirement for zero-point quantization. This modification streamlines both the software implementation and hardware-level optimization, and based on our empirical evaluations, leads to only minor differences in overall model effectiveness.

\paragraph{LLaMA-alike Components.}

LLaMA~\cite{llama, llama2} has become the standard architectural template for open-source large language models. Aligning with open-source trends, the design of \text{BitNet b1.58} adopts key elements of the LLaMA model family, such as RMSNorm~\cite{rmsnorm}, SwiGLU~\cite{swiglu}, rotary embedding~\cite{rope}, and the removal of bias terms. These architectural choices enable \text{BitNet b1.58} to be compatible with widely used open-source platforms (for example, Huggingface, vLLM~\cite{vllm}, and llama.cpp\footnote{\url{https://github.com/ggerganov/llama.cpp}}) with minimal integration effort.

\section{Results}

We conducted a head-to-head evaluation of \text{BitNet b1.58} against our FP16 \text{LLaMA LLM} reproduction across multiple model scales. Both sets of models underwent pre-training on the RedPajama dataset~\cite{redpajama}, processing 100 billion tokens to facilitate an equitable assessment. Their zero-shot capabilities were then benchmarked on a comprehensive suite of language understanding benchmarks, including ARC-Easy~\cite{arc}, ARC-Challenge~\cite{arc}, Hellaswag~\cite{hellaswag}, Winogrande~\cite{winoGrande}, PIQA~\cite{piqa}, OpenbookQA~\cite{openbookqa}, and BoolQ~\cite{boolq}. Additionally, we provided validation perplexity scores on the WikiText2~\cite{wikitext2} and C4~\cite{c4} corpora.

In addition to accuracy metrics, we analyzed both GPU memory utilization and inference latency for \text{LLaMA LLM} and \text{BitNet b1.58}. All latency figures were obtained using the FasterTransformer\footnote{\url{https://github.com/NVIDIA/FasterTransformer}} library, a high-performance codebase optimized for large language model inference on GPUs. \text{BitNet b1.58} utilized Ladder’s~\cite{ladder} 2-bit kernel for its experiments. We report inference costs in terms of time per generated token, reflecting the dominant factor in model deployment.

Table~\ref{tab:ppl} presents a comparative overview of perplexity and computational resource demands for both \text{BitNet b1.58} and \text{LLaMA LLM}. According to our experiments, \text{BitNet b1.58} achieves comparable perplexity to the full-precision \text{LLaMA LLM} beginning at the 3B scale, all while providing a 2.71$\times$ increase in speed and consuming only 28\% of the GPU memory. For the 3.9B configuration, \text{BitNet b1.58} further demonstrates a 2.4$\times$ speedup and a 3.32$\times$ reduction in memory, all the while substantially outperforming the 3B variant of \text{LLaMA LLM}.

Detailed zero-shot evaluation results for downstream tasks are found in Table~\ref{tab:zero-shot}. Evaluations were carried out following the standardized approach provided by \emph{lm-evaluation-harness}\footnote{\url{https://github.com/EleutherAI/lm-evaluation-harness}}. The data indicate that the performance differential between \text{BitNet b1.58} and \text{LLaMA LLM} narrows as model capacity grows. Importantly, at the 3B size, \text{BitNet b1.58} effectively closes the gap with the FP16 baseline. Echoing the perplexity trends, end-task evaluations confirm that the 3.9B \text{BitNet b1.58} not only exceeds the performance of the 3B \text{LLaMA LLM} but also does so with lower computational overhead. Collectively, these findings position \text{BitNet b1.58} as a Pareto-dominant alternative to leading large language model architectures.

\paragraph{Memory and Latency}

The model was further expanded to 7B, 13B, and 70B parameters, and we analyzed the associated costs. As depicted in Figure~\ref{fig:latency-memory-energy}, both latency and memory usage exhibit notable trends: speed-up improvements become more pronounced with increased model sizes. Notably, the \text{BitNet b1.58} model at 70B parameters achieves a 4.1-fold increase in speed relative to the \text{LLaMA LLM} baseline. This acceleration is primarily attributed to the growth in computation time for the \emph{nn.Linear} operation as model sizes increase. The memory utilization demonstrates a parallel pattern, given that embeddings are maintained in full precision and account for a reduced proportion of the overall memory footprint in larger models. These measurements of latency and memory were obtained utilizing a 2-bit kernel, indicating that further optimizations remain possible for additional cost reductions.

\paragraph{Energy}

In addition, we estimated the energy required for arithmetic operations in both \text{BitNet b1.58} and \text{LLaMA LLM}. Our analysis centers on matrix multiplication, as it represents the dominant computational cost within LLMs. Figure~\ref{fig:energy} breaks down the components of overall energy usage. In \text{BitNet b1.58}, the majority of computation involves INT8 additions, whereas \text{LLaMA LLM} relies on FP16 addition and multiplication. Drawing upon the energy estimation framework described in~\cite{energycost, pokebnn}, we found that \text{BitNet b1.58} offers a 71.4-fold reduction in arithmetic energy consumption for matrix multiplications on 7nm hardware. Furthermore, we recorded the total end-to-end energy required for models processing 512 tokens. Results indicate that as model size increases, \text{BitNet b1.58} continues to demonstrate significantly greater energy efficiency relative to the FP16 \text{LLaMA LLM} baseline. This efficiency gain can be attributed to the rising share of \emph{nn.Linear} operations as models grow, whereas other architectural components consume a relatively smaller fraction of the total cost in larger models.

\paragraph{Throughput}

We evaluate and contrast the throughput of \text{BitNet b1.58} and \text{LLaMA LLM} at the 70B parameter scale, leveraging two 80GB A100 GPUs. To enable \text{LLaMA LLM} 70B to fit on these devices, pipeline parallelism~\cite{gpipe} is employed. Batch size is systematically increased up to the available GPU memory capacity, using a fixed sequence length of 512. As indicated in Table~\ref{tab:throughput}, \text{BitNet b1.58} 70B can accommodate batch sizes up to 11 times larger than those supported by \text{LLaMA LLM}, culminating in an 8.9-fold increase in throughput.

\textbf{\text{BitNet b1.58} introduces a novel scaling law connecting inference cost with model performance.} To illustrate, using data from Figure~\ref{fig:latency-memory-energy} and \ref{fig:energy}, we can relate equivalent efficiency between different model sizes when using 1.58-bit precision as opposed to the conventional 16-bit:

\begin{itemize}
    \item BitNet b1.58 at 13B outperforms a 3B FP16 LLM in latency, memory footprint, and energy usage.
    \item BitNet b1.58 at 30B is more effective than a 7B FP16 LLM with respect to latency, memory, and energy requirements.
    \item BitNet b1.58 at 70B surpasses a 13B FP16 LLM on the same efficiency metrics.
\end{itemize}

\paragraph{Training with 2T Tokens}

Training token count is a pivotal variable in LLM effectiveness. To probe the scaling properties of \text{BitNet b1.58} with respect to token throughput, we trained a \text{BitNet b1.58} instance on 2T tokens, following the data pipeline of StableLM-3B~\cite{StableLM-3B-4E1T}, recognized as a leading open-source 3B model. Evaluation was conducted over a suite of benchmarks, namely Winogrande~\cite{winoGrande}, PIQA~\cite{piqa}, SciQ~\cite{sciq}, LAMBADA~\cite{lambada}, and ARC-easy~\cite{arc}, with zero-shot accuracy summarized in Table~\ref{tab:2t}. For tasks providing both accuracy and normalized accuracy, their average is reported. The performance numbers for StableLM 3b with 2T tokens are sourced from its official technical documentation. Our experiments demonstrate that \text{BitNet b1.58} delivers improved outcomes across all evaluated tasks, thereby highlighting the robust generalization capacity of 1.58-bit LLMs.

\section{Discussion and Future Work}

\textbf{1-bit Mixture-of-Experts (MoE) LLMs}

Mixture-of-Experts (MoE) methods have emerged as an efficient solution for scaling LLMs in terms of computational resources. Although MoE architectures significantly cut down floating point operations (FLOPs), their practical adoption is hampered by substantial memory demands and the associated overhead from inter-chip communication. 1.58-bit LLMs provide a pathway to overcome these limitations. Reducing memory requirements decreases the number of hardware units necessary for deploying MoE models, while simultaneously alleviating the cost of data transfer for activations between networked devices. In optimal scenarios where the full model fits on a single chip, such overhead can potentially be eliminated.

\textbf{Native Support of Long Sequence in LLMs}

Managing extended context lengths is increasingly vital in large-scale language modeling. A core barrier to efficient long sequence inference stems from the memory usage of KV caches. \text{BitNet b1.58} marks notable progress toward native long-context capabilities by shrinking activation storage from 16 bits to 8 bits, which makes it possible to double the context window within fixed memory constraints. For 1.58-bit LLMs, additional (lossless) compression down to 4 bits—or possibly less—could be achievable, a direction we propose to investigate in future research.

\textbf{LLMs on Edge and Mobile}

1.58-bit LLMs have considerable promise for enabling advanced language models on resource-constrained edge and mobile platforms. Such environments often struggle with tight memory and limited compute budgets, hampering the deployment of large models. The low memory and reduced power consumption characteristics of 1.58-bit LLMs unlock the potential for supporting sophisticated LLM-driven applications on these devices, opening new possibilities previously unattainable. Additionally, these quantized models are well-suited to CPU architectures prevalent in edge and mobile devices, which means that \text{BitNet b1.58} can be efficiently utilized, further advancing on-device language capabilities.

\textbf{New Hardware for 1-bit LLMs}

Recent advances such as Groq\footnote{\url{https://groq.com/}} have demonstrated the viability and promise of building specialized hardware (e.g., LPUs) designed for LLM workloads. Building upon these developments, we advocate for dedicated research and development aimed at new systems and hardware architectures tailored for 1-bit LLMs, taking advantage of the innovative computation paradigms introduced by BitNet~\cite{bitnet}.

\end{document}